\lecture{13}{2026-10-07}{More transfer theorems}

\begin{proposition}[{\cite[Proposition 2.17]{Melczer2021}}]
    If
    \[
        F(z) = (1-z)^{-\alpha} \left( \log \frac{1}{1-z} \right)^\beta
    \]
    where \(\alpha \in \mathbb{C} \setminus \mathbb{Z}_{\le 0}\) or \(\beta \in \mathbb{C} \setminus \mathbb{Z}_{\ge 0}\), then
    \[
        [z^n] F(z) = \frac{n^{\alpha - 1}}{\Gamma(\alpha)} \left( \log n \right)^\beta \left( 1 + \frac{c_1}{\log n} + \frac{c_2}{(\log n)^2} + \cdots \right)
    \]
    where
    \[
        c_k = (-1)^k \binom{\beta}{k} \Gamma(\alpha) \left. \frac{d^k}{ds^k} \frac{1}{\Gamma(s)} \right|_{s = \alpha}.
    \]
\end{proposition}

\begin{proof}[Proof idea]
    We write the Cauchy integral with the substitution \(z = 1 + t/n\):
    \[
        F(1 + t/n)(1 + t/n)^{-n-1} \approx e^{-t} (-t)^{-\alpha} (\log n)^\beta \left( 1 - \frac{\beta \log(-t)}{\log n} + \cdots \right),
    \]
    and evaluate along a Hankel contour \(\gamma\) using
    \begin{align*}
        \frac{1}{2\pi i} \int_{\gamma} e^{-t} (-t)^s \log^k(-t) \, dt
        &= (-1)^k \frac{d^k}{ds^k} \left( \frac{1}{2\pi i} \int_{\gamma} (-t)^s e^{-t} \, dt \right) \\
        &= (-1)^k \frac{d^k}{ds^k} \frac{1}{\Gamma(s)}.
    \end{align*}
\end{proof}

\begin{remark}[{\cite[Remark 2.4]{Melczer2021}}]
    If \(\alpha \in \mathbb{Z}_{\le 0}\) and \(\beta \in \mathbb{Z}_{\ge 0}\), the Gamma pole is avoided by differentiating \(\log \frac{1}{1-z}\) directly.
\end{remark}

\begin{example}
    \[
        \frac{d}{dz} \log\left(\frac{1}{1-z}\right) = \frac{1}{1-z} = \sum_{n \ge 0} z^n \implies [z^n] \log\left(\frac{1}{1-z}\right) = \frac{1}{n}.
    \]
\end{example}

\begin{example}
    \[
        [z^n] (1-z)^m \log\left(\frac{1}{1-z}\right) = (-1)^m \frac{m!}{n(n-1)\cdots(n-m)}.
    \]
\end{example}

\section*{Transfers}

If \(R > 0\) and \(\epsilon \in (0, \pi/2)\), then the \emph{\(\Delta\)-domain} is defined as
\[
    \Delta(R, \epsilon) = \{ z \in \mathbb{C} : |z| < R + \epsilon, \, z \neq R, \text{ and } |\operatorname{Arg}(z-R)| \ge \pi/2 - \epsilon \}.
\]

Since \([z^n] F(\rho z) = \rho^{n} [z^n] F(z)\), we can assume \(R = 1\) without loss of generality.

\begin{theorem}[Transfer theorem, {\cite[Proposition 2.17]{Melczer2021}, \cite[Theorem VI.3, Corollary VI.1]{FlajoletSedgewick2009}}]
    Let \(F(z) = \sum_{n \ge 0} f_n z^n\) be analytic in a \(\Delta\)-domain \(\Delta(1, \epsilon)\) for some \(\epsilon > 0\) and suppose that
    \[
        G(z) = (1-z)^{-\alpha} \left( \log \frac{1}{1-z} \right)^\beta = \sum_{n \ge 0} g_n z^n.
    \]
    Then, as \(n \to \infty\):
    \begin{itemize}
        \item if \(F(z) = O(G(z))\) as \(z \to 1\) in \(\Delta(1, \epsilon)\), then \(f_n = O(g_n) = O\left(n^{\alpha - 1} (\log n)^\beta\right)\);
        \item if \(F(z) = o(G(z))\) as \(z \to 1\) in \(\Delta(1, \epsilon)\), then \(f_n = o(g_n) = o\left(n^{\alpha - 1} (\log n)^\beta\right)\);
        \item if \(F(z) \sim G(z)\) as \(z \to 1\) in \(\Delta(1, \epsilon)\) and \(\alpha \notin \mathbb{Z}_{\le 0}\), then \(f_n \sim g_n \sim \frac{n^{\alpha - 1}}{\Gamma(\alpha)} (\log n)^\beta\).
    \end{itemize}
\end{theorem}

\begin{remark}
    The proof proceeds by deforming the Cauchy contour to the boundary of the \(\Delta\)-domain and evaluating the integrals along the rectilinear and circular segments.
\end{remark}

\begin{example}
    Consider the generating function for Catalan numbers:
    \[
        F(z) = \frac{1 - \sqrt{1 - 4z}}{2z} = \sum_{n \ge 0} C_n z^n.
    \]
    The dominant singularity is at \(z = 1/4\).
    Expanding locally around \(z = 1/4\) gives
    \[
        F(z) = 2 - 4\left(\frac{1}{4} - z\right)^{1/2} + 8\left(\frac{1}{4} - z\right) - 16\left(\frac{1}{4} - z\right)^{3/2} + O\left(\left(\frac{1}{4} - z\right)^2\right).
    \]
    Applying the transfer theorem term by term:
    \begin{align*}
        [z^n] \left(\frac{1}{4} - z\right)^{1/2} &= \frac{4^n}{2} \frac{n^{-3/2}}{\Gamma(-1/2)} \left( 1 + \frac{(-1/2)(-3/2)}{n} + O\left(\frac{1}{n^2}\right) \right), \\
        [z^n] \left(\frac{1}{4} - z\right)^{3/2} &= \frac{4^n}{2} \frac{n^{-5/2}}{\Gamma(-3/2)} \left( 1 + \frac{(-3/2)(-5/2)}{n} + O\left(\frac{1}{n^2}\right) \right), \\
        [z^n] O\left(\left(\frac{1}{4} - z\right)^2\right) &= O\left(4^n n^{-3}\right).
    \end{align*}
    Therefore,
    \[
        C_n = \frac{1}{n+1} \binom{2n}{n} = \frac{4^n}{\sqrt{\pi n^3}} \left( 1 - \frac{9}{8n} + O\left(\frac{1}{n^2}\right) \right).
    \]
\end{example}
